\documentclass[10pt,a4paper]{amsart}
\usepackage{amsmath,amssymb,amsthm,graphicx}
\usepackage{verbatim}
\newenvironment{reference}{\comment}{\endcomment}
%\newenvironment{reference}{}{}
\newenvironment{slogan}{\comment}{\endcomment}
\newenvironment{history}{\comment}{\endcomment}

% For commutative diagrams we use Xy-pic
\usepackage[all]{xy}

% We use 2cell for 2-commutative diagrams.
\xyoption{2cell}
% Theorem environments.
%
\theoremstyle{plain}
\newtheorem{theorem}[subsection]{Theorem}
\newtheorem{proposition}[subsection]{Proposition}
\newtheorem{lemma}[subsection]{Lemma}

\theoremstyle{definition}
\newtheorem{definition}[subsection]{Definition}
\newtheorem{example}[subsection]{Example}
\newtheorem{exercise}[subsection]{Exercise}
\newtheorem{situation}[subsection]{Situation}

\theoremstyle{remark}
\newtheorem{remark}[subsection]{Remark}
\newtheorem{remarks}[subsection]{Remarks}

\numberwithin{equation}{subsection}

% Macros
%
\def\lim{\mathop{\mathrm{lim}}\nolimits}
\def\colim{\mathop{\mathrm{colim}}\nolimits}
\def\Spec{\mathop{\mathrm{Spec}}}
\def\Hom{\mathop{\mathrm{Hom}}\nolimits}
\def\Ext{\mathop{\mathrm{Ext}}\nolimits}
\def\SheafHom{\mathop{\mathcal{H}\!\mathit{om}}\nolimits}
\def\SheafExt{\mathop{\mathcal{E}\!\mathit{xt}}\nolimits}
\def\Sch{\mathit{Sch}}
\def\Mor{\mathop{\mathrm{Mor}}\nolimits}
\def\Ob{\mathop{\mathrm{Ob}}\nolimits}
\def\Sh{\mathop{\mathit{Sh}}\nolimits}
\def\NL{\mathop{N\!L}\nolimits}
\def\CH{\mathop{\mathrm{CH}}\nolimits}
\def\proetale{{pro\text{-}\acute{e}tale}}
\def\etale{{\acute{e}tale}}
\def\QCoh{\mathit{QCoh}}
\def\Ker{\mathop{\mathrm{Ker}}}
\def\Im{\mathop{\mathrm{Im}}}
\def\Coker{\mathop{\mathrm{Coker}}}
\def\Coim{\mathop{\mathrm{Coim}}}

% Boxtimes
%
\DeclareMathSymbol{\boxtimes}{\mathbin}{AMSa}{"02}

%
% Macros for moduli stacks/spaces
%
\def\QCohstack{\mathcal{QC}\!\mathit{oh}}
\def\Cohstack{\mathcal{C}\!\mathit{oh}}
\def\Spacesstack{\mathcal{S}\!\mathit{paces}}
\def\Quotfunctor{\mathrm{Quot}}
\def\Hilbfunctor{\mathrm{Hilb}}
\def\Curvesstack{\mathcal{C}\!\mathit{urves}}
\def\Polarizedstack{\mathcal{P}\!\mathit{olarized}}
\def\Complexesstack{\mathcal{C}\!\mathit{omplexes}}
% \Pic is the operator that assigns to X its picard group, usage \Pic(X)
% \Picardstack_{X/B} denotes the Picard stack of X over B
% \Picardfunctor_{X/B} denotes the Picard functor of X over B
\def\Pic{\mathop{\mathrm{Pic}}\nolimits}
\def\Picardstack{\mathcal{P}\!\mathit{ic}}
\def\Picardfunctor{\mathrm{Pic}}
\def\Deformationcategory{\mathcal{D}\!\mathit{ef}}
\usepackage{cleveref}
\crefname{originallemma}{original lemma}{original lemmas}
\Crefname{originallemma}{Original lemma}{Original lemmas}
\makeatletter
\newlabel{more-algebra-lemma-injective-abelian}{{01D7}{source}{}{sourcepointer-1}{}}
\newlabel{more-algebra-lemma-injective-abelian@cref}{{[originallemma][0][]01D7}{[1][0][]source}{}{}{}}
\makeatother
\crefname{originalequation}{original equation}{original equations}
\Crefname{originalequation}{Original equation}{Original equations}
\makeatletter
\newlabel{equation-free-sheaf-on}{{05NY}{source}{}{sourcepointer-2}{}}
\newlabel{equation-free-sheaf-on@cref}{{[originalequation][0][]05NY}{[1][0][]source}{}{}{}}
\makeatother
\title{A criterion for injectivity of a sheaf of abelian groups}
\author{The Stacks Project Authors}
\date{stacks-project:injectives.tex:local-snapshot-20261012; GNU FDL 1.2 or later}
\begin{document}
\maketitle
\noindent\textit{Editorial note.} All delivered mathematical text is the author's own, copied from the pinned Stacks Project chapter file sites-cohomology.tex; no statement or proof step was rewritten, reordered or removed.\par\medskip
\noindent\textit{Editorial note.} The proof is the author's complete argument; the gluing step is displayed as a strictly bigger subsheaf obtained by adding an image, and no step is left to the reader.\par\medskip
\noindent\textit{Editorial note.} Every labelled reference occurring inside the counted spans is carried in assembly.external\_references as a non-load-bearing navigation pointer (same-chapter ones at their definition site, cross-chapter ones at the invocation inside the counted span), so no load-bearing mathematics is replaced by a stub. No bibliography citation occurs inside the counted spans.\par\medskip
\noindent\textit{Editorial note.} The same-name check for this card scanned every staged Stacks chapter of this repository -- all sources/stacks-*/pinned-original-source.tex files, the card's own chapter excluded -- and the labels found are recorded in dedup.same\_name\_cross\_chapter on the card.\par\medskip
\noindent\textit{Editorial note.} The preamble is the chapter's own retained minimal preamble (sources/stacks-sites-cohomology/retained-preamble.tex), a verbatim excerpt of the native preamble: the Xy-pic setup and the native theorem declarations with the mathematical macros.\par\medskip
\noindent\textit{Allowed standard prerequisites (editorial).}\begin{itemize}
\item The sheaf associated to a representable presheaf. Same chapter, the construction of the sheaf of abelian groups associated to the presheaf represented by an object and the discussion of its subsheaves. Supplies the standard objects and their subsheaves out of which the test morphisms in the criterion are built.
\item Gluing of morphisms on a sheaf. Same chapter, the sheaf axioms for morphisms of sheaves of abelian groups. Supplies that two morphisms agreeing on the intersection of their domains glue to a morphism on the sum of the subsheaves, which is the mechanism of the proof.
\item Injective sheaves of abelian groups. Same chapter, the definition of an injective object and the extension criterion used for sheaves. Supplies the extension property that the criterion has to verify.
\end{itemize}\medskip
\ifdefined\hypertarget\hypertarget{sourcepointer-1}{}\fi
\noindent\textit{Source reference.} Original lemma 01D7: sources/stacks-more-algebra/pinned-original-source.tex, source lines 13553--13553. Navigation pointer invoked at proof line 1031; the counted span contains the pointer only and the cited result is not reproduced inside it, so the pointer is registered as a non-load-bearing navigation stub in the established form for Stacks material. No bibliography citation is involved.\par\medskip
\ifdefined\hypertarget\hypertarget{sourcepointer-2}{}\fi
\noindent\textit{Source reference.} Original equation 05NY: sources/stacks-injectives/pinned-original-source.tex, source lines 1008--1008. Navigation pointer invoked at proof line 1042; the counted span contains the pointer only and the cited result is not reproduced inside it, so the pointer is registered as a non-load-bearing navigation stub in the established form for Stacks material. No bibliography citation is involved.\par\medskip
\begin{lemma}
\label{lemma-characterize-injectives}
Suppose $\mathcal{J}$ is a sheaf of abelian groups with the following
property: For all $X\in \Ob(\mathcal{C})$, for any abelian subsheaf
$\mathcal{S} \subset \mathbf{Z}_X^\#$ and any morphism
$\varphi : \mathcal{S} \to \mathcal{J}$, there exists a morphism
$\mathbf{Z}_X^\# \to \mathcal{J}$ extending $\varphi$.
Then $\mathcal{J}$ is an injective sheaf of abelian groups.
\end{lemma}

\begin{proof}
Let $\mathcal{F} \to \mathcal{G}$ be an injective map
of abelian sheaves. Suppose $\varphi : \mathcal{F} \to \mathcal{J}$
is a morphism. Arguing as in the proof of
More on Algebra, Lemma \ref{more-algebra-lemma-injective-abelian}
we see that it suffices
to prove that if $\mathcal{F} \not = \mathcal{G}$, then we
can find an abelian sheaf $\mathcal{F}'$,
$\mathcal{F} \subset \mathcal{F}' \subset \mathcal{G}$
such that (a) the inclusion $\mathcal{F} \subset \mathcal{F}'$ is strict,
and (b) $\varphi$ can be extended to $\mathcal{F}'$.
To find $\mathcal{F}'$, let $X$ be an object of $\mathcal{C}$ such
that the inclusion $\mathcal{F}(X) \subset \mathcal{G}(X)$
is strict. Pick $s \in \mathcal{G}(X)$, $s \not \in \mathcal{F}(X)$.
Let $\psi : \mathbf{Z}_X^\# \to \mathcal{G}$ be the morphism corresponding
to the section $s$ via (\ref{equation-free-sheaf-on}). Set
$\mathcal{S} = \psi^{-1}(\mathcal{F})$. By assumption the morphism
$$
\mathcal{S} \xrightarrow{\psi} \mathcal{F} \xrightarrow{\varphi} \mathcal{J}
$$
can be extended to a morphism $\varphi' : \mathbf{Z}_X^\# \to \mathcal{J}$.
Note that $\varphi'$ annihilates the kernel of $\psi$ (as this is true
for $\varphi$). Thus $\varphi'$ gives rise to a morphism
$\varphi'' : \Im(\psi) \to \mathcal{J}$
which agrees with $\varphi$ on the intersection
$\mathcal{F} \cap \Im(\psi)$ by construction.
Thus $\varphi$ and $\varphi''$ glue to give an extension
of $\varphi$ to the strictly bigger subsheaf
$\mathcal{F}' = \mathcal{F} + \Im(\psi)$.
\end{proof}

\end{document}
